\documentclass{amsart}
% For dealing with references we use the comment environment
\usepackage{verbatim}
\newenvironment{reference}{\comment}{\endcomment}
%\newenvironment{reference}{}{}
\newenvironment{slogan}{\comment}{\endcomment}
\newenvironment{history}{\comment}{\endcomment}

% For commutative diagrams we use Xy-pic
\usepackage[all]{xy}

% We use 2cell for 2-commutative diagrams.
\xyoption{2cell}
\UseAllTwocells

% We use multicol for the list of chapters between chapters
\usepackage{multicol}

% This is generally recommended for better output
\usepackage{lmodern}
\usepackage[T1]{fontenc}

% For cross-file-references

% Package for hypertext links:
\usepackage{hyperref}

% For any local file, say "hello.tex" you want to link to please

% Theorem environments.
%
\theoremstyle{plain}
\newtheorem{theorem}[subsection]{Theorem}
\newtheorem{proposition}[subsection]{Proposition}
\newtheorem{lemma}[subsection]{Lemma}

\theoremstyle{definition}
\newtheorem{definition}[subsection]{Definition}
\newtheorem{example}[subsection]{Example}
\newtheorem{exercise}[subsection]{Exercise}
\newtheorem{situation}[subsection]{Situation}

\theoremstyle{remark}
\newtheorem{remark}[subsection]{Remark}
\newtheorem{remarks}[subsection]{Remarks}

\numberwithin{equation}{subsection}

% Macros
%
\def\lim{\mathop{\mathrm{lim}}\nolimits}
\def\colim{\mathop{\mathrm{colim}}\nolimits}
\def\Spec{\mathop{\mathrm{Spec}}}
\def\Hom{\mathop{\mathrm{Hom}}\nolimits}
\def\Ext{\mathop{\mathrm{Ext}}\nolimits}
\def\SheafHom{\mathop{\mathcal{H}\!\mathit{om}}\nolimits}
\def\SheafExt{\mathop{\mathcal{E}\!\mathit{xt}}\nolimits}
\def\Sch{\mathit{Sch}}
\def\Mor{\mathop{\mathrm{Mor}}\nolimits}
\def\Ob{\mathop{\mathrm{Ob}}\nolimits}
\def\Sh{\mathop{\mathit{Sh}}\nolimits}
\def\NL{\mathop{N\!L}\nolimits}
\def\CH{\mathop{\mathrm{CH}}\nolimits}
\def\proetale{{pro\text{-}\acute{e}tale}}
\def\etale{{\acute{e}tale}}
\def\QCoh{\mathit{QCoh}}
\def\Ker{\mathop{\mathrm{Ker}}}
\def\Im{\mathop{\mathrm{Im}}}
\def\Coker{\mathop{\mathrm{Coker}}}
\def\Coim{\mathop{\mathrm{Coim}}}

% Boxtimes
%
\DeclareMathSymbol{\boxtimes}{\mathbin}{AMSa}{"02}

%
% Macros for moduli stacks/spaces
%
\def\QCohstack{\mathcal{QC}\!\mathit{oh}}
\def\Cohstack{\mathcal{C}\!\mathit{oh}}
\def\Spacesstack{\mathcal{S}\!\mathit{paces}}
\def\Quotfunctor{\mathrm{Quot}}
\def\Hilbfunctor{\mathrm{Hilb}}
\def\Curvesstack{\mathcal{C}\!\mathit{urves}}
\def\Polarizedstack{\mathcal{P}\!\mathit{olarized}}
\def\Complexesstack{\mathcal{C}\!\mathit{omplexes}}
% \Pic is the operator that assigns to X its picard group, usage \Pic(X)
% \Picardstack_{X/B} denotes the Picard stack of X over B
% \Picardfunctor_{X/B} denotes the Picard functor of X over B
\def\Pic{\mathop{\mathrm{Pic}}\nolimits}
\def\Picardstack{\mathcal{P}\!\mathit{ic}}
\def\Picardfunctor{\mathrm{Pic}}
\def\Deformationcategory{\mathcal{D}\!\mathit{ef}}

\setlength{\emergencystretch}{3em}
\begin{document}
\title[Stacks Project, Tag 0FD7]{589 --- Sufficiently negative ample twists have no global sections for coherent sheaves without closed associated points on proper schemes}
\maketitle
\noindent Copyright (C) 2005--2025 Johan de Jong. Stacks Project authors. Source: \href{https://stacks.math.columbia.edu/tag/0FD7}{Tag 0FD7}.

\noindent Editorial difficulty note (not author text). Difficulty H2: The proof must control extensions of arbitrary coherent sheaves rather than only integral structure sheaves. It proves an inductive support property by globally generating annihilator ideals of zero-dimensional pieces, then rules out negative sections on positive-dimensional integral supports by asymptotic Riemann\textendash{}Roch..

\noindent Editorial provenance note (not author text). History: excerpt from revision \texttt{a0444\allowbreak 6e57e\allowbreak c1fbc\allowbreak 252a8\allowbreak 71afc\allowbreak ec775\allowbreak 2fb28\allowbreak 07b14}, varieties.tex, lines 11074--11155. Reference presentation was adapted; original-author context and core support blocks were retained or added. The mathematical wording is unchanged. Licensed under GNU FDL 1.2 or later; no invariant sections or cover texts. See ../licenses/COPYING. Context blocks reproduce author definitions and supporting results; they are not additional counted problems.

\begin{lemma}
\label{lemma-vanishin-h0-negative}
Let $k$ be a field. Let $X$ be a proper scheme over $k$. Let $\mathcal{L}$
be an ample invertible $\mathcal{O}_X$-module. Let $\mathcal{F}$ be a
coherent $\mathcal{O}_X$-module. If $\text{Ass}(\mathcal{F})$ does not
contain any closed points, then
$\Gamma(X, \mathcal{F} \otimes_{\mathcal{O}_X} \mathcal{L}^{\otimes n}) = 0$
for $n \ll 0$.
\end{lemma}

\begin{proof}
For a coherent $\mathcal{O}_X$-module $\mathcal{F}$
let $\mathcal{P}(\mathcal{F})$ be the property: there
exists an $n_0 \in \mathbf{Z}$ such that for $n \leq n_0$
every section $s$ of
$\mathcal{F} \otimes_{\mathcal{O}_X} \mathcal{L}^{\otimes n}$
has support consisting only of closed points.
Since $\text{Ass}(\mathcal{F}) =
\text{Ass}(\mathcal{F} \otimes_{\mathcal{O}_X} \mathcal{L}^{\otimes n})$
we see that it suffices to prove $\mathcal{P}$
holds for all coherent modules on $X$.
To do this we will prove that conditions (1), (2), and (3) of
Cohomology of Schemes, Lemma
\href{https://stacks.math.columbia.edu/tag/01YM}{lemma property higher rank cohomological (Tag 01YM)}
are satisfied.

\medskip\noindent
To see condition (1) suppose that
$$
0 \to \mathcal{F}_1 \to \mathcal{F} \to \mathcal{F}_2 \to 0
$$
is a short exact sequence of coherent $\mathcal{O}_X$-modules
such that we have $\mathcal{P}$ for $\mathcal{F}_i$, $i = 1, 2$.
Let $n_1, n_2$ be the cutoffs we find.
Let $\mathcal{F}'_2 \subset \mathcal{F}_2$ be the maximal
coherent submodule whose support is a finite set of closed
points. Let $\mathcal{I} \subset \mathcal{O}_X$ be the annihilator
of $\mathcal{F}'_2$. Since $\mathcal{L}$ is ample, we can find an $e > 0$
such that $\mathcal{I} \otimes_{\mathcal{O}_X} \mathcal{L}^{\otimes e}$
is globally generated. Set $n_0 = \min(n_2, n_1 - e)$. Let $n \leq n_0$
and let $t$ be a global section of
$\mathcal{F} \otimes \mathcal{L}^{\otimes n}$.
The image of $t$ in $\mathcal{F}_2 \otimes \mathcal{L}^{\otimes n}$
falls into $\mathcal{F}'_2 \otimes \mathcal{L}^{\otimes n}$ because
$n \leq n_2$. Hence for any
$s \in \Gamma(X, \mathcal{I} \otimes_{\mathcal{O}_X} \mathcal{L}^{\otimes e})$
the product $t \otimes s$ lies in
$\mathcal{F}_1 \otimes \mathcal{L}^{\otimes n + e}$.
Thus $t \otimes s$ has support contained in
the finite set of closed points in $\text{Ass}(\mathcal{F}_1)$
because $n + e \leq n_1$. Since by our choice of $e$
we may choose $s$ invertible in any point
not in the support of $\mathcal{F}'_2$
we conclude that the support of $t$ is contained in the union of the 
finite set of closed points in $\text{Ass}(\mathcal{F}_1)$ and
the finite set of closed points in $\text{Ass}(\mathcal{F}_2)$.
This finishes the proof of condition (1).

\medskip\noindent
Condition (2) is immediate.

\medskip\noindent
For condition (3) we choose $\mathcal{G} = \mathcal{O}_Z$.
In this case, if $Z$ is a closed point of $X$, then there is
nothing the show. If $\dim(Z) > 0$, then we will show that
$\Gamma(Z, \mathcal{L}^{\otimes n}|_Z) = 0$
for $n < 0$. Namely, let $s$ be a nonzero section of a negative
power of $\mathcal{L}|_Z$. Choose a nonzero section $t$ of a positive
power of $\mathcal{L}|_Z$ (this is possible as $\mathcal{L}$ is ample, see
Properties, Proposition \href{https://stacks.math.columbia.edu/tag/01Q3}{proposition characterize ample (Tag 01Q3)}).
Then $s^{\deg(t)} \otimes t^{\deg(s)}$
is a nonzero global section of $\mathcal{O}_Z$
(because $Z$ is integral) and hence
a unit (Lemma \href{https://stacks.math.columbia.edu/tag/0BUG}{lemma proper geometrically reduced global sections (Tag 0BUG)}).
This implies that $t$ is a trivializing section
of a positive power of $\mathcal{L}$.
Thus the function $n \mapsto \dim_k \Gamma(X, \mathcal{L}^{\otimes n})$
is bounded on an infinite set of positive integers which contradicts
asymptotic Riemann-Roch
(Proposition \href{https://stacks.math.columbia.edu/tag/0BJ8}{proposition asymptotic riemann roch (Tag 0BJ8)})
since $\dim(Z) > 0$.
\end{proof}
\end{document}
